% !TEX root = ../../thesis.tex
\chapter{Tổng quan}

Chương này giới thiệu hai bài toán mà khóa luận giải quyết ---
tóm tắt tin tức tự động và kiểm chứng tin tức tự động --- kỹ thuật
Mixture-of-Agents (MoA) hỗ trợ đóng góp kỹ thuật, và khảo sát
các công trình liên quan trong nghiên cứu tiếng Việt cũng như quốc
tế. Chương kết thúc bằng phát biểu về khoảng trống mà khóa luận
lấp đầy.

\section{Bài toán tóm tắt văn bản tự động}

\textbf{Định nghĩa.}
Tóm tắt văn bản tự động là bài toán tạo ra một phiên bản ngắn hơn
của văn bản nguồn mà vẫn bảo toàn thông tin thiết yếu. Cho một tài
liệu đầu vào $D$, bộ tóm tắt $f$ tạo ra bản tóm tắt $S = f(D)$
sao cho $|S| \ll |D|$ và $S$ giữ lại các dữ kiện then chốt của $D$.

\textbf{Phân loại.}
Các hệ thống tóm tắt thường được phân loại theo hai trục.

\begin{itemize}
  \item \textit{Trích xuất so với trừu tượng hóa.}
        Hệ thống trích xuất chọn lựa các câu (hoặc đoạn câu) hiện có
        từ văn bản gốc. Hệ thống trừu tượng hóa tạo ra văn bản mới
        và có thể diễn giải lại, sắp xếp lại thứ tự hoặc kết hợp
        các ý tưởng. Các bộ tóm tắt dựa trên LLM hiện đại chủ yếu
        là kiểu trừu tượng hóa.
  \item \textit{Đơn tài liệu so với đa tài liệu.}
        Hệ thống đơn tài liệu tóm tắt một bài báo riêng lẻ; hệ thống
        đa tài liệu tổng hợp thông tin từ nhiều bài báo liên quan. Hệ
        thống trong khóa luận này là đơn tài liệu, nhưng pipeline MoA
        có thể được hiểu như một bộ tổng hợp đa nguồn suy biến, trong
        đó ``nhiều nguồn'' chính là các bản thảo LLM khác nhau của cùng
        một bài báo.
\end{itemize}

\textbf{Thách thức đặc thù của tiếng Việt.}
Tiếng Việt đặt ra ba thách thức cho bài toán tóm tắt.
\begin{enumerate}
  \item \textit{Phân đoạn từ (Word segmentation).} Tiếng Việt được viết
        với khoảng trắng giữa các âm tiết, không phải giữa các từ; một từ
        khái niệm có thể kéo dài hai đến ba âm tiết
        (\textit{tóm tắt}, \textit{kiểm chứng}). Các hệ số overlap từ vựng
        được tính trên đơn vị token khoảng trắng sẽ hệ thống hóa việc
        không ghi nhận công cho các cách diễn đạt lại tái cấu trúc thuật
        ngữ đa âm tiết.
  \item \textit{Dấu phụ và thanh điệu.} Tiếng Việt có sáu thanh từ vựng,
        được mã hóa dưới dạng dấu phụ trên nguyên âm; kho từ vựng của bộ
        mã hóa phân mảnh các ký tự có dấu sẽ mất đi độ chi tiết ngữ nghĩa.
  \item \textit{Khan hiếm corpus đánh giá.} So với tiếng Anh
        (CNN/DailyMail, XSum, MultiNews), tiếng Việt không có bộ đánh giá
        tóm tắt được chú thích bởi con người được chấp nhận rộng rãi với
        các bản tóm tắt tham chiếu. Đây là lý do thực tế khiến ROUGE/BLEU/
        BERTScore trong nghiên cứu tiếng Việt thường được tính
        \textit{so với bài báo gốc}, không phải so với bản tóm tắt tham
        chiếu của con người --- một lưu ý phương pháp luận sẽ được đề cập
        lại ở Chương~4.
\end{enumerate}

\section{Bài toán kiểm chứng tin tức}

\textbf{Tin sai lệch trong báo chí tiếng Việt.}
Người đọc Việt Nam thường xuyên gặp phải các tiêu đề giật gân, trích
dẫn phiến diện và thông tin sai lệch lan truyền trên mạng xã hội và
các tổng hợp tin tức. Kiểm chứng thủ công không đủ khả năng đáp ứng
quy mô thông tin hiện tại; một công cụ tự động bổ sung, có khả năng
gắn cờ các tuyên bố không có căn cứ, ngày càng trở nên có giá trị.

\textbf{Pipeline tiêu chuẩn.}
Tài liệu học thuật hội tụ về một pipeline kiểm chứng tự động ba bước.
\begin{enumerate}
  \item \textit{Phát hiện tuyên bố (Claim detection).} Trích xuất các
        tuyên bố nguyên tử đáng kiểm chứng từ tài liệu.
  \item \textit{Truy xuất bằng chứng (Evidence retrieval).} Với mỗi
        tuyên bố, truy xuất các đoạn văn từ kho ngữ liệu (hoặc web mở)
        có thể hỗ trợ hoặc bác bỏ tuyên bố đó.
  \item \textit{Phán quyết (Verdict prediction).} Phân loại mỗi tuyên
        bố là \textit{được xác nhận (entailed)},
        \textit{bị bác bỏ (contradicted)} hoặc \textit{không đề cập
        (not mentioned)} dựa trên bằng chứng, và tổng hợp các phán quyết
        cấp tuyên bố thành điểm tin cậy cấp tài liệu.
\end{enumerate}

\textbf{Xác minh tăng cường tìm kiếm.}
Tìm kiếm web mở đã trở thành nguồn bằng chứng trung thực nhất cho
kiểm chứng, vì không có cơ sở tri thức nào có thể theo kịp tin tức
tức thì. Xác minh tăng cường tìm kiếm kết hợp một API tìm kiếm web
với một LLM đọc các đoạn trích xuất được và đưa ra phán quyết; đây
là thiết kế được áp dụng trong Chương~3 (Tavily $\rightarrow$ GPT-4o-mini).

\section{Phương pháp Mixture-of-Agents (MoA)}

\textbf{Nguồn gốc.}
Wang et al.\ đề xuất Mixture-of-Agents vào năm 2024 \cite{Wang2024}.
Quan sát thực nghiệm trung tâm là một LLM đóng vai trò \textit{tác
nhân tổng hợp} cho ra câu trả lời tốt hơn khi được cung cấp nhiều
câu trả lời ứng viên từ các LLM khác --- kể cả khi các LLM kia riêng
lẻ yếu hơn chính tác nhân tổng hợp. Các tác giả gọi đây là
\textit{khả năng hợp tác của các LLM}.

\textbf{Kiến trúc.}
MoA là một cấu trúc phân tầng. Tầng $\ell$ gồm $n$ tác nhân
$\{A_{\ell,1}, \ldots, A_{\ell,n}\}$; mỗi tác nhân ở tầng $\ell+1$
nhận lời nhắc cùng với đầu ra của tất cả các tác nhân ở tầng $\ell$.
Về mặt hình thức, tác nhân $i$ ở tầng $\ell+1$ tạo ra

\[
y_{\ell+1,i} \;=\; A_{\ell+1,i}\!\left(\,[y_{\ell,1}, y_{\ell,2},
\ldots, y_{\ell,n}]\,+\,x\right),
\]

\noindent trong đó $x$ là lời nhắc của người dùng và $+$ biểu thị
phép nối chuỗi. Trong phiên bản hai tầng (được sử dụng trong khóa
luận này), tầng~1 là tầng \textit{tác nhân đề xuất} và tầng~2 là một
\textit{tác nhân tổng hợp} duy nhất.

\textbf{Kết quả đã công bố.}
Trên bộ đánh giá AlpacaEval~2.0, một cấu hình MoA với sáu tác nhân
đề xuất mã nguồn mở và một tác nhân tổng hợp đạt tỷ lệ thắng LC(length-controlled) là
$65,1\%$, vượt qua GPT-4o ($57,5\%$) theo cùng giao thức đánh giá.
Hiệu suất tăng đơn điệu theo số lượng tầng và số lượng tác nhân đề
xuất mỗi tầng (Bảng~3--4 của \cite{Wang2024}).

\textbf{Câu hỏi mở cho khóa luận.}
Bài báo MoA đánh giá trên các bộ đánh giá tuân thủ hướng dẫn
(AlpacaEval, MT-Bench, FLASK) trong đó LLM-giám khảo --- không phải
độ overlap với bất kỳ văn bản nguồn nào --- là tiêu chuẩn đánh giá.
Liệu MoA có chuyển giao sang \textit{tóm tắt có neo ngữ cảnh gốc}
(nơi tồn tại một bài báo nguồn và có thể đo được độ overlap), và liệu
lợi ích có xuất hiện trong tiếng Việt hay không, chính là câu hỏi
thực nghiệm mà khóa luận này trả lời.

\section{Các công trình liên quan}

\textbf{Hệ thống tóm tắt tiếng Việt.}
\textit{Google News Việt Nam} và \textit{Báo Mới} cung cấp tin tức
tổng hợp với các đoạn mô tả ngắn, nhưng các đoạn đó là trích dẫn trực
tiếp từ đoạn dẫn đầu bài, không phải bản tóm tắt biên tập trừu tượng.
Nghiên cứu học thuật về tóm tắt tiếng Việt đã sử dụng ViT5
\cite{Nguyen2022} cho việc tạo văn bản trừu tượng hóa, được đánh giá
gần như độc quyền bằng ROUGE so với bản tóm tắt tham chiếu (khi có)
hoặc so với bài báo gốc (phổ biến hơn).

\textbf{Công cụ kiểm chứng thông tin.}
\textit{ClaimBuster} (UT Arlington) phát hiện các tuyên bố đáng kiểm
chứng trong bài phát biểu chính trị tiếng Anh; \textit{Google Fact
Check Tools} API tổng hợp các phán quyết từ các tổ chức kiểm chứng
thông tin hiện có. Cả hai đều không nhắm vào tiếng Việt, và cả hai
đều không hoạt động ở cấp bài báo cho người đọc thông thường.

\textbf{Tiện ích mở rộng trình duyệt hỗ trợ đọc.}
Các tiện ích \textit{TLDR This}, \textit{Summari} và
\textit{Reader Mode} nhúng bản tóm tắt trong trang trên các trang web
tùy ý nhưng ưu tiên tiếng Anh; bộ tóm tắt của chúng hoạt động kém
với tiếng Việt, và chúng không cung cấp tính năng kiểm chứng thông
tin. \textit{NewsGuard} và \textit{Bias Buddy} đánh giá độ tin cậy
của \textit{nhà xuất bản}, không phải bài báo riêng lẻ, và dựa trên
đánh giá thủ công thay vì xác minh từng bài báo.

\textbf{Đánh giá bằng LLM-giám khảo.}
Zheng et al.\ \cite{Zheng2023} giới thiệu MT-Bench và phương pháp
Chatbot Arena, chứng minh rằng một LLM mạnh (GPT-4) đồng ý với sở
thích của con người đối với đầu ra chatbot ở mức $\sim$80\% (xấp xỉ
mức mà hai người đánh giá đồng ý với nhau). FLASK \cite{Ye2023} mở
rộng phương pháp này sang đánh giá rubric chi tiết theo nhiều kỹ năng
phù hợp. AlpacaEval \cite{Dubois2024} bổ sung ngẫu nhiên hóa vị trí
và kiểm soát độ dài để giảm thiểu hai thiên kiến đã biết của LLM-giám
khảo. Module đánh giá Trục~B của khóa luận này được xây dựng trên các
phương pháp này.

\textbf{Khoảng trống.}
Theo hiểu biết của tôi, chưa có công trình nào trước đây (i)~áp dụng
Mixture-of-Agents cho tiếng Việt, (ii)~kết hợp phân phối sản phẩm dựa
trên tiện ích mở rộng với kiểm chứng thông tin và tóm tắt trên tin tức
tiếng Việt, hoặc (iii)~tường minh đối chiếu đánh giá overlap, LLM-giám
khảo và đánh giá con người trong một khung \textit{ba trục} duy nhất
cho bài toán tóm tắt tiếng Việt. Khóa luận này sẽ nhắm tới việc khai thác và lấp đầy 
ba khoảng trống đó.

\section{Tổng kết chương}

Chương này đã đặt ra bài toán tóm tắt và kiểm chứng tin tức tiếng Việt
như hai nhu cầu song song mà khóa luận giải quyết, giới thiệu
Mixture-of-Agents như kỹ thuật mà khóa luận áp dụng, khảo sát các
công trình liên quan trong nghiên cứu tiếng Việt và quốc tế, và xác
định khoảng trống: chưa có hệ thống tiếng Việt nào kết hợp dung hợp
MoA, triển khai qua tiện ích mở rộng trình duyệt, và khung đánh giá ba
trục về mặt phương pháp luận. Chương~2 phát triển nền tảng lý thuyết
mà hệ thống trong Chương~3 được xây dựng dựa trên đó.
